% Источники: exam/materials/преподаватель/2020/23-03-2020-MODEL L5(1).doc; exam/materials/моделирование.pdf, разделы 31 и 33; GitHub HumsterProgrammer/iu9-notes/8sem/Моделирование/Моделирование_лекция-07_05-03-26.md; /Users/egor/Desktop/сем8/Моделирование/рк/РкМод/варианты/вариант1.pdf; /Users/egor/Desktop/сем8/Моделирование/рк/РкМод/варианты/вариант2.pdf.
\section{Метод моментов и метод максимального правдоподобия.}

\subsection*{Метод моментов (ММ)}

\textbf{Теоретические моменты} распределения с параметром $\theta$:
\[
  \mu_r(\theta) = \int_{-\infty}^{+\infty} x^r\,f(x,\theta)\,dx, \quad r = 1, 2, \ldots
\]

\textbf{Выборочные моменты}:
\[
  \hat{\mu}_r = \frac{1}{n}\sum_{i=1}^n x_i^r.
\]

\textbf{Алгоритм}:
\begin{enumerate}
  \item Выразить теоретические моменты $\mu_1(\theta),\ldots,\mu_k(\theta)$ через неизвестные параметры.
  \item Приравнять теоретические моменты выборочным: $\mu_r(\theta) = \hat{\mu}_r$.
  \item Решить систему уравнений относительно $\theta$.
\end{enumerate}

\textbf{Пример (ненаблюдаемые одновременно, $\xi_2 = a\xi_1^{b}$).}
Логарифмируем: $\ln\xi_2 = \ln a + b\ln\xi_1$. По раздельным выборкам:
\[
  \overline{\ln X} = \frac{1}{n}\sum \ln x_i, \quad
  s^2_{\ln X} = \frac{1}{n-1}\sum(\ln x_i - \overline{\ln X})^2,
\]
аналогично $\overline{\ln Y}$, $s^2_{\ln Y}$. Оценки:
\[
  \hat b = \sqrt{\frac{s^2_{\ln Y}}{s^2_{\ln X}}}, \qquad
  \hat a = \exp\!\bigl(\overline{\ln Y} - \hat b\,\overline{\ln X}\bigr).
\]

\subsection*{Пример из варианта 1/2: модифицированный метод моментов}

Пусть требуется построить формулы для случая линейной функциональной зависимости между ненаблюдаемыми одновременно характеристиками:
\[
  \xi_2 = a\xi_1 + b.
\]
По раздельным выборкам оцениваются средние $\overline{x}$, $\overline{y}$ и выборочные дисперсии $s_x^2$, $s_y^2$. Для линейной зависимости:
\[
  \overline{y} = a\overline{x} + b, \qquad
  s_y^2 = a^2s_x^2.
\]
Следовательно,
\[
  \hat a = \sqrt{\frac{s_y^2}{s_x^2}}, \qquad
  \hat b = \overline{y} - \hat a\,\overline{x}.
\]
Положительный корень берётся при предполагаемой прямой зависимости. Если по смыслу предметной области связь обратная, знак коэффициента $a$ задаётся предметной областью, а модуль коэффициента определяется из отношения дисперсий.

\subsection*{Метод максимального правдоподобия (ММП)}

\textbf{Функция правдоподобия}:
\[
  L(\theta) = \prod_{i=1}^n f(x_i, \theta).
\]

\textbf{Логарифмическое правдоподобие}:
\[
  \ln L(\theta) = \sum_{i=1}^n \ln f(x_i, \theta).
\]

\textbf{Уравнение правдоподобия} (при одном параметре):
\[
  \frac{\partial \ln L}{\partial \theta} = 0.
\]

\textbf{Алгоритм}:
\begin{enumerate}
  \item Записать функцию правдоподобия $L(\theta)$.
  \item Перейти к логарифму $\ln L(\theta)$.
  \item Дифференцировать по каждому параметру, приравнять к нулю.
  \item Решить систему уравнений правдоподобия.
\end{enumerate}

\textbf{Пример из основного материала: логистическая регрессия.}
Параметры модели выбираются так, чтобы функция правдоподобия наблюдаемой выборки обращалась в максимум.

\clearpage
